% ============================================================================
% Apéndice: Recursos computacionales del proyecto CUEVA
% Insertar en 07_appendices.tex
% ============================================================================

\section{Recursos computacionales}
\label{ap:computo}

\primerparrafo Todo el cómputo del proyecto CUEVA se realizó en una única estación de trabajo de acceso dedicado, sin recurrir a infraestructura de clúster distribuida. Las especificaciones de la máquina y el volumen total de trabajo acumulado se documentan aquí como referencia para la reproducibilidad de los resultados.

% -----------------------------------------------------------------------
\subsection{Especificaciones de la máquina}
\label{ap:computo:maquina}

\primerparrafo La estación de trabajo utilizada es un nodo de cómputo con las siguientes características (Tabla~\ref{ap:tab:maquina}):

\begin{table}[H]\centering\small
\caption{Especificaciones del nodo de cómputo del proyecto CUEVA.}
\label{ap:tab:maquina}
\begin{tabular}{ll}
\hline
\textbf{Componente}   & \textbf{Especificación} \\
\hline
Procesador        & AMD Ryzen Threadripper 3990X (Zen~2) \\
Núcleos / hilos   & 64 núcleos físicos / 128 hilos \\
Frecuencia        & 2.9~GHz base $\rightarrow$ 4.35~GHz turbo \\
Arquitectura NUMA & 1 nodo NUMA \\
Memoria RAM       & 251~GiB (256~GB) DDR4 \\
Almacenamiento    & SSD NVMe de 1.8~TB \\
Paralelismo       & OpenMP (hasta 64 hilos por corrida) \\
\hline
\end{tabular}
\end{table}

% -----------------------------------------------------------------------
\subsection{Trabajo realizado y tiempo de cómputo}
\label{ap:computo:trabajo}

\primerparrafo El proyecto se desarrolló de forma continua desde julio de 2024 hasta junio de 2026 (cerca de dos años), y comprendió tres grandes bloques de simulaciones de la inestabilidad de Kelvin--Helmholtz en magnetohidrodinámica relativista resistiva, más el post-procesamiento. El costo computacional de cada bloque se resume en la Tabla~\ref{ap:tab:computo}.

\begin{table}[H]\centering\small
\caption{Bloques de simulaciones y costo computacional acumulado.}
\label{ap:tab:computo}
\begin{tabular}{p{6.8cm} c c r}
\hline
\textbf{Bloque de simulaciones} & \textbf{Corridas} & \textbf{Resolución} & \textbf{Core-h} \\
\hline
Campañas magnéticas A/B/C ($\sigma=6000,\,10\,000$)$^{\,\dagger}$  & 119 & $512\times256$  &  19\,910       \\
Barrido de conductividad $\sigma\in[100,\,40\,000]$$^{\,\ddagger}$  &  52 & $512\times256$  & $\sim$8\,000   \\
Simulación de alta resolución$^{\,\ddagger}$                         &   1 & $1024\times512$ & $\sim$1\,800   \\
Post-procesamiento$^{\,\ddagger}$                                    & --- & ---             & $\sim$100      \\
\hline
\textbf{Total estimado}                                              &     &                 & $\sim$\textbf{29\,800} \\
\hline
\end{tabular}
\begin{minipage}{0.90\linewidth}
\vspace{4pt}\footnotesize
$^{\dagger}$~Medido directamente con \texttt{/usr/bin/time\,-v}; 60 de las 119 ejecuciones terminaron
en crash numérico, parte del estudio de estabilidad del campo magnético.\\[2pt]
$^{\ddagger}$~Estimado a partir de la resolución, el número de pasos temporales y el costo medido de
las corridas equivalentes de las campañas A/B/C.
\end{minipage}
\end{table}

\primerparrafo Las $\sim$29\,800 core-horas acumuladas equivalen a $\sim$1\,240 core-días
($\approx3.4$~años-núcleo), o bien a unos 19.5~días de los 64 núcleos trabajando sin interrupción.
El período de mayor intensidad computacional abarcó de febrero a junio de 2026 (barrido de
conductividad y campañas magnéticas).
La memoria utilizada alcanzó un pico de $\sim$9~GB por corrida, con un máximo de $\sim$36~GB
al ejecutar cuatro simulaciones en paralelo; el volumen total de datos generados fue de
$\sim$522~GB en disco.

Sin el acceso a esta máquina habría sido inviable acumular las cerca de 30\,000 horas-núcleo
que sostienen este trabajo. Se expresa sincero agradecimiento a quienes facilitaron su uso,
recurso sin el cual ni el barrido de conductividad, ni el estudio de variación del campo
magnético, ni las simulaciones de alta resolución hubieran sido posibles.